% !TEX program = lualatex
% =====================================================================
% Defensa de tesis, DECK v3, estilo Metropolis (modo oscuro).
% Universidad Tecnologica de Pereira · MISC · 2026.
% Autores: Alejandro Gomez Huertas & Juan Diego Garzon Ovalle.
% Compilar con LuaLaTeX (dos veces):  lualatex beamer_defensa_v3.tex
% =====================================================================
\documentclass[aspectratio=169,11pt]{beamer}
\usepackage{fontspec}
\usepackage[spanish,es-noshorthands]{babel}
\usepackage{xcolor}
\usepackage{graphicx}
\usepackage{booktabs}
\usepackage{tikz}
\usetikzlibrary{arrows.meta}
\graphicspath{{../tesis_latex/}{fig/}}

\usetheme{metropolis}
\metroset{sectionpage=progressbar, progressbar=frametitle, numbering=fraction, block=fill}

% -------------------------------------------------- paleta (modo oscuro)
\definecolor{darkbg}{rgb}{0.05,0.10,0.125}
\definecolor{card}{rgb}{0.10,0.15,0.185}
\definecolor{greytx}{rgb}{0.87,0.88,0.90}
\definecolor{amber}{rgb}{1.0,0.80,0.25}
\definecolor{orange}{rgb}{1.0,0.55,0.40}
\definecolor{teal}{rgb}{0.40,0.80,0.92}
\definecolor{seafoam}{rgb}{0.45,0.86,0.70}
\definecolor{coral}{rgb}{1.0,0.46,0.42}
\definecolor{muted}{rgb}{0.62,0.68,0.78}
\definecolor{navy}{rgb}{0.62,0.80,1.0}   % remapeado a azul claro (visible en oscuro)

% -------------------------------------------------- colores del tema
\setbeamercolor{normal text}{fg=greytx,bg=darkbg}
\setbeamercolor{background canvas}{bg=darkbg}
\setbeamercolor{alerted text}{fg=amber}
\setbeamercolor{example text}{fg=seafoam}
\setbeamercolor{frametitle}{fg=amber,bg=darkbg}
\setbeamercolor{progress bar}{fg=amber,bg=card}
\setbeamercolor{title separator}{fg=amber}
\setbeamercolor{progress bar in section page}{fg=amber,bg=card}
\setbeamercolor{itemize item}{fg=amber}
\setbeamercolor{itemize subitem}{fg=teal}
\setbeamercolor{enumerate item}{fg=amber}
\setbeamercolor{standout}{fg=amber,bg=darkbg}
\setbeamercolor{block title}{fg=darkbg,bg=amber}
\setbeamercolor{block body}{fg=greytx,bg=card}
\setbeamercolor{block title alerted}{fg=white,bg=coral}
\setbeamercolor{block body alerted}{fg=greytx,bg=card}
\setbeamercolor{block title example}{fg=darkbg,bg=seafoam}
\setbeamercolor{block body example}{fg=greytx,bg=card}

% -------------------------------------------------- helper: figura sobre tarjeta blanca
\newcommand{\figcard}[2][\linewidth]{\tikz\node[fill=white,rounded corners=2pt,inner sep=2.5pt]{\includegraphics[width=#1]{#2}};}

% -------------------------------------------------- captions numerados (protocolo de tesis)
\newcounter{figc}
\newcounter{tabc}
\newcommand{\figcap}[1]{\stepcounter{figc}\par\vspace{2pt}{\scriptsize\color{muted}\textbf{Figura~\thefigc.}~#1}}
\newcommand{\tabcap}[1]{\stepcounter{tabc}{\scriptsize\color{muted}\textbf{Tabla~\thetabc.}~#1\par}\vspace{3pt}}

% -------------------------------------------------- datos del titulo
\title{Estabilidad de Métodos de Explicabilidad (XAI) en GNNs}
\subtitle{Detección de Lavado de Dinero bajo Desbalance Extremo}
\author{Alejandro Gómez Huertas \and Juan Diego Garzón Ovalle}
\institute{{\color{amber}\textsc{Universidad Tecnológica de Pereira}}\\ Facultad de Ingenierías \textperiodcentered\ Maestría en Ingeniería en Sistemas y Computación\\ Director: Ph.D. Cristian Rosero Arias}
\date{Trabajo de grado para optar al título de Magíster \textperiodcentered\ Pereira, Colombia \textperiodcentered\ 2026}

\begin{document}

\maketitle

\begin{frame}{Contenido de la sustentación}
\begin{enumerate}
  \item[\textcolor{amber}{\textbf{1}}] \textbf{Problema y motivación} \hfill {\small\color{muted}detección de lavado y la opacidad de las GNN}
  \item[\textcolor{amber}{\textbf{2}}] \textbf{Marco conceptual} \hfill {\small\color{muted}arquitecturas, explicadores y las tres propiedades}
  \item[\textcolor{amber}{\textbf{3}}] \textbf{Metodología: dos ejes} \hfill {\small\color{muted}Elliptic real + grafo sintético con \emph{ground-truth}}
  \item[\textcolor{amber}{\textbf{4}}] \textbf{Resultados} \hfill {\small\color{muted}estabilidad, plausibilidad, fidelidad}
  \item[\textcolor{amber}{\textbf{5}}] \textbf{Conclusiones} \hfill {\small\color{muted}aportes, límites y trabajo futuro}
\end{enumerate}
\vspace{1em}
\begin{block}{Presentación a dos voces}
\small \textbf{Alejandro}: contexto y metodología. \quad \textbf{Juan Diego}: resultados y cierre. \quad Duración objetivo: 33 a 36 min.
\end{block}
\end{frame}

% ##################################################################### 1
\section{Problema y motivación}

\begin{frame}{El lavado de dinero es un problema de red}
\begin{columns}[T]
\begin{column}{0.54\textwidth}
\textbf{\color{navy}Tres fases clásicas}
\begin{enumerate}\small
  \item \textbf{Colocación}: el dinero ilícito entra al sistema.
  \item \textbf{Estratificación}: se fragmenta en muchas transacciones para borrar el rastro.
  \item \textbf{Integración}: vuelve a la economía con apariencia legal.
\end{enumerate}
\vspace{0.4em}
\textbf{\color{navy}Tipologías = patrones entre cuentas}
\begin{itemize}\small
  \item \emph{Structuring} (pitufeo), \emph{layering}, \emph{fan-in}, \emph{fan-out}.
\end{itemize}
\end{column}
\begin{column}{0.46\textwidth}
\begin{block}{La idea clave}
Una transacción \textbf{aislada} parece normal. Lo que delata el lavado es el \alert{patrón de conexiones} entre cuentas.
\end{block}
\vspace{0.3em}
{\footnotesize\color{muted} Por eso se modela como \textbf{grafo}, no como una tabla de transacciones independientes.}
\end{column}
\end{columns}
\end{frame}

\begin{frame}{El problema: detectar sin caja negra}
\begin{columns}[T]
\begin{column}{0.56\textwidth}
\begin{itemize}
  \item El lavado mueve \alert{entre el 2 y el 5\% del PIB mundial}.
  \item Los sistemas \textbf{por reglas} generan \alert{entre el 95 y el 98\% de falsos positivos}: reglas rígidas que los criminales aprenden a esquivar.
  \item Las \textbf{GNN} modelan la red de flujos y \textbf{superan a lo tabular}, que mira cada transacción aislada.
  \item \textbf{Pero son cajas negras}: no dicen \emph{por qué} marcan una transacción.
\end{itemize}
\vspace{0.3em}
{\scriptsize\color{muted} Cifras: informes UNODC y GAFI/FATF (ver Cap. 2 del manuscrito).}
\end{column}
\begin{column}{0.44\textwidth}\centering
\figcard{chapter_2/images_ch2/enfoques_moeny_laundering.png}
\figcap{Enfoques de detección: reglas frente a modelos de grafos. \emph{Fuente: elaboración propia.}}
\end{column}
\end{columns}
\end{frame}

\begin{frame}{Por qué la caja negra es inaceptable aquí}
\begin{itemize}
  \item Una alerta puede \alert{congelar cuentas} o iniciar una investigación: el costo de un error es alto.
  \item En un entorno \textbf{regulado}, el analista debe \textbf{justificar cada decisión} ante un auditor.
  \item Y debe ser \textbf{reproducible}: la misma entrada tiene que dar la misma explicación.
\end{itemize}
\vspace{0.5em}
\begin{block}{Consecuencia}
La explicabilidad no es un lujo, es un \alert{requisito}. Y si la explicación \textbf{cambia} cada vez que se recalcula, no sirve para auditar: por eso la \alert{estabilidad} es la propiedad crítica y el foco de esta tesis.
\end{block}
\end{frame}

\begin{frame}{Pregunta de investigación y objetivos}
\begin{block}{Pregunta}
¿Cómo se comporta la \alert{estabilidad} de los métodos XAI sobre GNNs para detección de lavado bajo desbalance, y qué combinación de \textbf{arquitectura}, \textbf{explicador} y \textbf{balanceo} produce la interpretación más robusta y auditable?
\end{block}
\vspace{0.4em}
\begin{columns}[T]
\begin{column}{0.5\textwidth}
\begin{itemize}\small
  \item[\textcolor{amber}{\textbf{O1}}] ¿Se degrada la estabilidad al agravarse el desbalance?
  \item[\textcolor{amber}{\textbf{O2}}] ¿Qué tan resilientes son GCN, GraphSAGE, GAT y TAGCN?
\end{itemize}
\end{column}
\begin{column}{0.5\textwidth}
\begin{itemize}\small
  \item[\textcolor{amber}{\textbf{O3}}] ¿Influyen las estrategias de balanceo?
  \item[\textcolor{amber}{\textbf{O4}}] Condensar todo en una matriz de recomendación.
\end{itemize}
\end{column}
\end{columns}
\end{frame}

\begin{frame}{Tres hipótesis falsables}
Nos comprometimos con tres predicciones \textbf{antes} de ver los datos:
\vspace{0.5em}
\begin{itemize}
  \item[\textcolor{coral}{\textbf{H1}}] La estabilidad \textbf{se degradará} al agravarse el desbalance.
  \item[\textcolor{coral}{\textbf{H2}}] \textbf{TAGCN} será la más estable, por su alcance multi-salto.
  \item[\textcolor{coral}{\textbf{H3}}] Un explicador más \textbf{estable} señalará mejor el \textbf{patrón real} (plausibilidad).
\end{itemize}
\vspace{0.6em}
{\footnotesize\color{muted} Comprometerse con predicciones concretas es lo que hace que refutarlas signifique algo. El desenlace de las tres, en Resultados.}
\end{frame}

\begin{frame}{La brecha en el estado del arte}
\begin{columns}[T]
\begin{column}{0.5\textwidth}
\begin{block}{Lo que ya se sabía}
\begin{itemize}\small
  \item Las GNN superan a lo tabular (Weber y colaboradores, 2019).
  \item TAGCN + SHAP con alto rendimiento (He y colaboradores, 2026).
  \item Predicción y explicabilidad, por separado.
\end{itemize}
\end{block}
\end{column}
\begin{column}{0.5\textwidth}
\begin{alertblock}{Lo que faltaba}
\begin{itemize}\small
  \item La \textbf{estabilidad} de las explicaciones sobre grafos AML no se había evaluado sistemáticamente.
  \item Elliptic no tiene \emph{ground-truth} de tipología: la plausibilidad no era medible.
\end{itemize}
\end{alertblock}
\end{column}
\end{columns}
\end{frame}

% ##################################################################### 2
\section{Marco conceptual}

\begin{frame}{Cómo funciona una GNN: paso de mensajes}
\begin{columns}[T]
\begin{column}{0.56\textwidth}
\begin{itemize}
  \item Cada nodo \textbf{actualiza su representación} agregando la de sus \textbf{vecinos}.
  \item Se repite \textbf{capa por capa}: con $L$ capas, un nodo "ve" hasta \textbf{$L$ saltos} de distancia (su \emph{campo receptivo}).
  \item La predicción sobre una transacción depende de \textbf{su vecindario}, no solo de ella.
\end{itemize}
\vspace{0.2em}
\begin{block}{Por eso, explicar =}
señalar \textbf{qué} del vecindario (qué aristas, qué \emph{features}) sostuvo la predicción.
\end{block}
\end{column}
\begin{column}{0.44\textwidth}\centering
\begin{tikzpicture}[>=Stealth, n/.style={circle,draw=teal,fill=card,text=greytx,minimum size=7mm,font=\scriptsize,inner sep=0}]
  \node[n,draw=amber,line width=1pt] (c) at (0,0) {$v$};
  \node[n] (a) at (-1.8,1.0) {};
  \node[n] (b) at (-2.0,-0.6) {};
  \node[n] (d) at (1.8,1.0) {};
  \node[n] (e) at (2.0,-0.6) {};
  \draw[->,teal,line width=1pt] (a)--(c);
  \draw[->,teal,line width=1pt] (b)--(c);
  \draw[->,teal,line width=1pt] (d)--(c);
  \draw[->,teal,line width=1pt] (e)--(c);
\end{tikzpicture}
\\[4pt]
{\scriptsize\color{muted} Los vecinos "envían mensajes" al nodo $v$.}
\end{column}
\end{columns}
\end{frame}

\begin{frame}{Cuatro arquitecturas GNN}
\begin{itemize}
  \item \textbf{\color{navy}GCN}: \emph{promedia} a los vecinos (convolución espectral de un salto). Simple y sólida.
  \item \textbf{\color{navy}GraphSAGE}: \emph{muestrea} vecinos y los agrega; inductiva, escala a grafos grandes.
  \item \textbf{\color{navy}GAT}: pondera a los vecinos con \emph{atención} --- no todos pesan igual.
  \item \textbf{\color{navy}TAGCN}: filtros polinómicos de orden $K$ que alcanzan \emph{varios saltos} de una vez.
\end{itemize}
\vspace{0.5em}
{\footnotesize\color{muted} Se eligieron por cubrir las familias dominantes de agregación de vecinos. Comparar las cuatro es el Objetivo 2.}
\end{frame}

\begin{frame}{Tres explicadores post-hoc}
{\footnotesize \textbf{Post-hoc} = explican un modelo \textbf{ya entrenado}, sin modificarlo. Cada uno produce una \textbf{máscara de aristas} y un \textbf{ranking de \emph{features}}.}
\vspace{0.5em}
\begin{itemize}
  \item \textbf{\color{navy}GNNExplainer}: optimiza una \emph{máscara por instancia} --- aprende qué importa para \emph{ese} nodo.
  \item \textbf{\color{navy}PGExplainer}: entrena una red \emph{amortizada} que aprende a explicar en todo el grafo (generaliza entre nodos).
  \item \textbf{\color{navy}GNNShap}: reparte el crédito con \emph{valores de Shapley} (teoría de juegos); muy consistente.
\end{itemize}
\vspace{0.4em}
{\footnotesize\color{muted} Cuatro arquitecturas $\times$ tres explicadores es el núcleo de la comparación.}
\end{frame}

\begin{frame}{Tres propiedades que no son lo mismo}
\begin{columns}[T]
\begin{column}{0.33\textwidth}
\begin{block}{Estabilidad}\small ¿La explicación se reproduce entre semillas y perturbaciones?\end{block}
\end{column}
\begin{column}{0.33\textwidth}
\begin{block}{Plausibilidad}\small ¿Señala el patrón real de lavado (\emph{ground-truth})?\end{block}
\end{column}
\begin{column}{0.33\textwidth}
\begin{block}{Fidelidad}\small ¿Refleja de verdad la decisión interna del modelo?\end{block}
\end{column}
\end{columns}
\vspace{0.8em}
\begin{center}\large
\alert{Tesis central:} \emph{las tres son empíricamente independientes.}
\end{center}
\vspace{0.3em}
{\footnotesize\color{muted} Un explicador puede ser \textbf{estable y equivocado} --- como un reloj parado: siempre marca lo mismo, y siempre mal.}
\end{frame}

% ##################################################################### 3
\section{Metodología}

\begin{frame}{Qué produce, en concreto, un explicador}
\begin{columns}[T]
\begin{column}{0.46\textwidth}
\centering
\begin{tikzpicture}[scale=0.80, every node/.style={transform shape}]
  \tikzset{
    nodo/.style={circle,draw=muted,fill=white,minimum size=13pt,inner sep=0pt,line width=0.6pt},
    foco/.style={circle,draw=coral,fill=coral!35,minimum size=16pt,inner sep=0pt,line width=1.1pt},
    rel/.style={draw=amber,line width=2.1pt},
    irr/.style={draw=muted!60,line width=0.7pt}
  }
  \node[foco] (c) at (0,0) {};
  \node[nodo] (a) at (-1.55,0.95) {};
  \node[nodo] (b) at (-1.75,-0.55) {};
  \node[nodo] (d) at (1.55,0.9) {};
  \node[nodo] (e) at (1.7,-0.65) {};
  \node[nodo] (f) at (0.15,-1.6) {};
  \node[nodo] (g) at (-2.9,0.35) {};
  \node[nodo] (h) at (2.85,0.1) {};
  \draw[rel] (c) -- (a);
  \draw[rel] (c) -- (d);
  \draw[rel] (a) -- (g);
  \draw[irr] (c) -- (b);
  \draw[irr] (c) -- (e);
  \draw[irr] (c) -- (f);
  \draw[irr] (d) -- (h);
  \draw[irr] (b) -- (f);
\end{tikzpicture}
\\[3pt]
{\footnotesize \textcolor{coral}{$\bullet$} nodo a explicar \quad \textcolor{amber}{\rule[2pt]{9pt}{2pt}} aristas relevantes}
\figcap{Esquema de una explicación: nodo objetivo, aristas relevantes y ranking de \emph{features}. \emph{Fuente: elaboración propia.}}
\end{column}
\begin{column}{0.52\textwidth}
\footnotesize
Ante una transacción marcada como ilícita, el explicador devuelve dos cosas:
\begin{enumerate}
  \item Una \textbf{máscara de aristas}: qué conexiones del vecindario sostienen la predicción.
  \item Un \textbf{ranking de \emph{features}}: qué atributos del nodo pesaron más, de los 166 disponibles.
\end{enumerate}
\begin{block}{Las tres preguntas}
\textbf{Estabilidad}: ¿sale lo mismo al repetir?\\
\textbf{Plausibilidad}: ¿coincide con el patrón real?\\
\textbf{Fidelidad}: ¿es lo que el modelo usó?
\end{block}
{\color{muted}En Elliptic el vecindario tiene una mediana de $\sim$2 nodos: ahí solo el ranking de \emph{features} discrimina.}
\end{column}
\end{columns}
\end{frame}

\begin{frame}{Metodología: un diseño de dos ejes}
\begin{columns}[T]
\begin{column}{0.5\textwidth}
\begin{block}{Eje 1 \textperiodcentered\ Elliptic (real)}
\begin{itemize}\small
  \item Validez \textbf{externa}: transacciones reales.
  \item Sin \emph{ground-truth}: solo estabilidad.
  \item Campo receptivo minúsculo ($\sim$2 nodos).
\end{itemize}
\end{block}
\end{column}
\begin{column}{0.5\textwidth}
\begin{block}{Eje 2 \textperiodcentered\ Sintético (propio)}
\begin{itemize}\small
  \item Validez \textbf{interna}: \emph{ground-truth} por nodo y arista.
  \item Permite medir plausibilidad y fidelidad.
  \item 4 tipologías estándar y 3 realizaciones del grafo.
\end{itemize}
\end{block}
\end{column}
\end{columns}
\end{frame}

\begin{frame}{Diseño factorial}
\begin{center}
{\color{teal}\Large\textbf{4}} arquitecturas \; $\times$ \; {\color{amber}\Large\textbf{3}} explicadores \; $\times$ \; {\color{seafoam}\Large\textbf{3}} balanceos \; $\times$ \; {\color{coral}\Large\textbf{5}} escenarios
\end{center}
\vspace{0.3em}
\begin{center}\small = 60 configuraciones por eje \, $\times$ 3 explicadores para el estudio de estabilidad\end{center}
\vspace{0.6em}
\begin{block}{Protocolo de robustez}
\small Estabilidad: 5 réplicas estocásticas por explicación. Inferencia (eje sintético): 3 semillas de modelo $\times$ 3 grafos, con Kruskal-Wallis, Wilcoxon e IC \emph{bootstrap}.
\end{block}
\end{frame}

\begin{frame}{Eje 1: Elliptic Bitcoin Dataset}
\begin{columns}[T]
\begin{column}{0.54\textwidth}
\begin{itemize}\small
  \item 203.769 nodos \textperiodcentered\ 234.355 aristas \textperiodcentered\ 166 features \textperiodcentered\ 49 pasos.
  \item Ilícitas $\sim$2,2\% (razón $\sim$1:9 en el subconjunto etiquetado).
  \item Partición temporal causal: entrenamiento en los pasos 1 a 34, validación de 35 a 42 y test de 43 a 49.
  \item Campo receptivo mediana $\sim$2 nodos: solo el Spearman discrimina.
\end{itemize}
\vspace{0.2em}
{\scriptsize\color{muted} Dataset: Weber et al. (2019).}
\end{column}
\begin{column}{0.46\textwidth}\centering
\figcard{chapter_4/images_ch4/dispersion_elliptic.png}
\figcap{Dispersión de la topología de Elliptic (grado y campo receptivo). \emph{Fuente: elaboración propia.}}
\end{column}
\end{columns}
\end{frame}

\begin{frame}{Eje 2: grafo sintético con ground-truth}
\begin{columns}[T]
\begin{column}{0.54\textwidth}
\begin{itemize}\small
  \item \textbf{¿Por qué?} medir plausibilidad exige saber cuál subgrafo es el patrón: Elliptic no lo da.
  \item 4 tipologías: \emph{structuring}, \emph{layering}, \emph{fan-in}, \emph{fan-out}.
  \item Aristas distractoras + firma atenuada: la métrica discrimina.
  \item \emph{Ground-truth} ciego al explicador y 3 realizaciones del grafo.
\end{itemize}
\end{column}
\begin{column}{0.46\textwidth}\centering
\figcard{chapter_5/images_ch5/hallazgo_estrella.png}
\figcap{Grafo sintético con tipologías de lavado y \emph{ground-truth} por arista. \emph{Fuente: elaboración propia.}}
\end{column}
\end{columns}
\end{frame}

\begin{frame}{Cómo se construyó el grafo sintético}
\begin{columns}[T]
\begin{column}{0.5\textwidth}
\footnotesize
\textbf{\color{navy}Decisiones que endurecen la prueba}
\begin{itemize}
  \item \textbf{Simetrización} de aristas: sin ella el campo receptivo cae a $\sim$2 nodos y la plausibilidad de subgrafo no es medible.
  \item \textbf{Aristas distractoras} ($n=3$ por nodo de patrón): sin ellas cualquier \emph{top-k} acertaría.
  \item \textbf{Firma de \emph{features} atenuada} de $+4$ a $+1{,}5$: evita que la tarea se resuelva sola.
\end{itemize}
{\color{muted}El \emph{ground-truth} se fija en la construcción, \textbf{antes} de correr ningún explicador.}
\end{column}
\begin{column}{0.5\textwidth}
\footnotesize
\textbf{\color{navy}Escala y replicación}
\begin{itemize}
  \item $\approx$9.500 nodos, $\approx$1.530 ilícitos, $\approx$31.000 aristas.
  \item 4 tipologías canónicas plantadas dentro.
  \item Grafo g0: factorial completo $\times$ 3 semillas (540 filas).
  \item Grafos g1 y g2: corte de verificación (144 filas).
\end{itemize}
{\color{muted}El contraste principal se calcula sobre la unión de los tres grafos.}
\end{column}
\end{columns}
\end{frame}

\begin{frame}{Métricas y protocolo estadístico}
\begin{itemize}
  \item \textbf{\color{teal}Estabilidad}: correlación de Spearman entre rankings de features (primaria).
  \item \textbf{\color{seafoam}Plausibilidad}: coincidencia con el \emph{ground-truth} de la tipología.
  \item \textbf{\color{amber}Fidelidad}: caída de la predicción al retirar lo importante (Fidelity+).
  \item \textbf{\color{coral}Rendimiento}: PR-AUC como métrica primaria (el F1 en umbral es degenerado bajo desbalance).
\end{itemize}
\vspace{0.6em}
\begin{block}{}
\small Estadística: Kruskal-Wallis, Wilcoxon e intervalos de confianza \emph{bootstrap} (pruebas no paramétricas: no asumen normalidad).
\end{block}
\end{frame}

% ##################################################################### 4
\section{Resultados}

\begin{frame}{Contribución: dos artefactos de evaluación}
\begin{columns}[T]
\begin{column}{0.5\textwidth}
\begin{alertblock}{1 \textperiodcentered\ Fallo de memoria silencioso}
\small Cálculo sobre el grafo completo: 13 configuraciones de GAT fallaban en silencio. \emph{Favorecía a GAT.}
\end{alertblock}
\end{column}
\begin{column}{0.5\textwidth}
\begin{alertblock}{2 \textperiodcentered\ Truncamiento del Spearman}
\small La métrica descartaba features fuera del umbral: rankings mutilados. \emph{Favorecía a GraphSAGE.}
\end{alertblock}
\end{column}
\end{columns}
\vspace{0.6em}
\begin{center}\small
\alert{Cada uno bastaba para una conclusión comparativa falsa.} La estabilidad depende del protocolo de medición.\\ Además: dos \emph{bugs} reportables en el PGExplainer de PyG 2.7.
\end{center}
\vspace{0.2em}
{\scriptsize\color{muted}\hfill En línea con la advertencia de Kosan et al. (2023) sobre la sensibilidad al protocolo.\hspace*{0.4cm}}
\end{frame}

\begin{frame}{El segundo artefacto: el truncamiento de la métrica}
\begin{columns}[T]
\begin{column}{0.52\textwidth}
\footnotesize
La métrica dimensionaba el vector de rangos por \texttt{top\_k} en vez de por el número real de \emph{features}.\\[4pt]
Con \texttt{top\_k}$=20$ sobre \textbf{166 features}, toda \emph{feature} con índice $\geq 20$ se descartaba: sobrevivían 2 o 3 y el resto quedaba empatado en cero.\\[4pt]
El truncamiento \textbf{no penaliza por igual}: castiga más a las arquitecturas cuya importancia se reparte sobre muchas \emph{features}.\\[4pt]
{\color{muted}Por eso deprimía a GAT ($+0{,}24$ al corregir) y a TAGCN ($+0{,}28$) mucho más que a GraphSAGE ($+0{,}10$), y fabricaba un liderazgo falso.}
\end{column}
\begin{column}{0.46\textwidth}
\begin{block}{Los dos artefactos, juntos}
\footnotesize
1. \textbf{Memoria}: 13 configs de GAT fallaban por OOM silencioso y su media se calculaba solo sobre las que terminaban. \emph{Favorecía a GAT.}\\[4pt]
2. \textbf{Métrica}: el truncamiento anterior. \emph{Favorecía a GraphSAGE.}\\[4pt]
Cada uno \textbf{por separado} bastaba para una conclusión falsa.
\end{block}
\end{column}
\end{columns}
\end{frame}

\begin{frame}{Resultado 1: la estabilidad separa dos grupos, no cuatro puestos}
\begin{columns}[T]
\begin{column}{0.5\textwidth}
\begin{itemize}\small
  \item[\textcolor{seafoam}{$\blacksquare$}] \textbf{GAT} \hfill $0{,}782 \pm 0{,}013$
  \item[\textcolor{seafoam}{$\blacksquare$}] \textbf{GCN} \hfill $0{,}758 \pm 0{,}025$
\end{itemize}
\vspace{2pt}
{\footnotesize\color{coral}\hspace{6pt}\textbf{brecha significativa entre grupos}}
\vspace{2pt}
\begin{itemize}\small
  \item[\textcolor{muted}{$\blacksquare$}] GraphSAGE \hfill $0{,}735 \pm 0{,}022$
  \item[\textcolor{muted}{$\blacksquare$}] TAGCN \hfill $0{,}672 \pm 0{,}077$
\end{itemize}
\vspace{0.5em}
\small Media de \textbf{3 semillas} (180 modelos). El $\pm$ es la \textbf{desviación típica entre semillas}, no el intervalo de confianza de la figura. Significativo \textbf{entre} grupos, \textbf{no dentro}: Wilcoxon pareado GAT vs GCN $p=0{,}17$ y GraphSAGE vs TAGCN $p=0{,}10$.
\end{column}
\begin{column}{0.5\textwidth}\centering
\figcard{ranking_khop_eb.png}
\figcap{Estabilidad (Spearman de GNNExplainer) por arquitectura en Elliptic, con 3 semillas, 180 modelos e IC 95\%. \emph{Fuente: elaboración propia.}}
\end{column}
\end{columns}
\end{frame}

\begin{frame}{Robustez: qué sostiene esa partición}
\begin{columns}[T]
\begin{column}{0.54\textwidth}
\footnotesize
\tabcap{Contrastes no paramétricos de la partición (Elliptic, 180 modelos).}
\begin{tabular}{lcc}
\toprule
\textbf{Contraste} & \textbf{Estadístico} & \textbf{$p$} \\
\midrule
Global (4 arq.) & Kruskal-Wallis & $2{,}8\!\times\!10^{-5}$ \\
Grupo alto vs bajo & Mann-Whitney & $1{,}4\!\times\!10^{-6}$ \\
\midrule
Dentro grupo alto & Kruskal-Wallis & $0{,}19$ \\
Dentro grupo bajo & Kruskal-Wallis & $0{,}08$ \\
\bottomrule
\end{tabular}
\vspace{0.5em}
\begin{itemize}\footnotesize
  \item Tamaño de efecto $\eta^2 = 0{,}13$ (grande), el mismo estadístico que reporta la tesis en el eje sintético.
  \item IC95\,\% bootstrap: se solapan \textbf{dentro}, apenas se tocan \textbf{entre} grupos.
  \item Sobrevive la corrección de \textbf{Holm} por comparaciones múltiples.
  \item Los contrastes \emph{dentro} de grupo son no pareados. Los $p$ de la lámina anterior ($0{,}17$ y $0{,}10$) son el Wilcoxon \textbf{pareado} sobre celdas comunes. Ambas vías no rechazan.
\end{itemize}
\end{column}
\begin{column}{0.44\textwidth}
\begin{block}{Reproducibilidad precisa}
\footnotesize
Reentrenamos la matriz completa. Los \textbf{pesos no} se reproducen bit a bit (\emph{scatter} atómicos en GPU).\\[4pt]
Las \textbf{conclusiones sí}: 25/60 sobre el filtro de calidad frente a 23/60, y la \textbf{misma partición}.
\end{block}
{\footnotesize\color{muted} Distinguir entre ambas también es una contribución metodológica.}
\end{column}
\end{columns}
\end{frame}

\begin{frame}{Resultado 1b: concordancia entre regímenes}
\begin{columns}[T]
\begin{column}{0.5\textwidth}
\begin{itemize}\small
  \item La \emph{inversión por densidad} era un \textbf{artefacto} de la métrica.
  \item Corregida, Elliptic (disperso) \textbf{concuerda} con el sintético (denso).
  \item El eje sintético separa \textbf{los mismos dos grupos}, sobre datos y \emph{pipeline} independientes.
\end{itemize}
\vspace{0.4em}
\begin{center}
\small Correlación de rangos entre regímenes\\[3pt]
{\color{coral}\Large\textbf{$-0{,}20$}} \; $\rightarrow$ \; {\color{amber}\Large\textbf{$+0{,}80$}}\\[2pt]
{\footnotesize métrica con bug \; $\rightarrow$ \; métrica corregida}\\[4pt]
{\footnotesize\color{muted} Sintético: GCN 0,97 y GAT 0,96 vs GraphSAGE 0,89 y TAGCN 0,89}
\end{center}
\end{column}
\begin{column}{0.5\textwidth}\centering
\figcard{chapter_5/images_ch5/contraste_regimen.png}
\figcap{Orden de estabilidad: Elliptic (disperso) frente al sintético (denso). \emph{Fuente: elaboración propia.}}
\end{column}
\end{columns}
\end{frame}

\begin{frame}{Resultado 2: disociación plausibilidad / fidelidad}
\begin{columns}[T]
\begin{column}{0.52\textwidth}
\begin{itemize}\small
  \item \textbf{Por un lado}, \textbf{\color{teal}PGExplainer} domina la plausibilidad: 0,80 vs 0,50, con 0,40 de azar (Wilcoxon $p\approx 2{,}6\times10^{-35}$).
  \item \textbf{Por otro}, \textbf{\color{coral}colapsa en fidelidad}: 0,11 vs 0,56 de GNNExplainer.
  \item \emph{El explicador más ``plausible'' no es el más fiel.}
\end{itemize}
\vspace{0.3em}
{\scriptsize\color{muted} La figura compara los dos explicadores que producen máscara de aristas. GNNShap no la produce por diseño: su plausibilidad es de \emph{features}, con otra escala y otra línea base, y va en la lámina de respaldo.}
\end{column}
\begin{column}{0.48\textwidth}\centering
\figcard{disociacion_eb.png}
\figcap{Plausibilidad de aristas frente a Fidelity+ (eje sintético, IC 95\% y azar de plausibilidad de aristas $=0{,}40$). \emph{Fuente: elaboración propia.}}
\end{column}
\end{columns}
\end{frame}

\begin{frame}{Resultado 3: el puente que no existe}
\begin{columns}[T]
\begin{column}{0.52\textwidth}
\begin{itemize}\small
  \item Nuestra hipótesis central: una explicación más estable señalaría mejor el patrón real.
  \item El dato dice que \textbf{no}. Correlación entre estabilidad y plausibilidad $r=-0{,}01$, IC 95\,\% de $-0{,}038$ a $+0{,}011$.
  \item El intervalo \textbf{incluye el cero}: no hay relación.
\end{itemize}
\vspace{0.4em}
\small\emph{Es un resultado nulo, y lo reportamos porque contradice lo que esperábamos.}
\end{column}
\begin{column}{0.48\textwidth}\centering
\figcard{chapter_5/images_ch5/puente_nulo.png}
\figcap{Relación entre estabilidad y plausibilidad (GNNExplainer): $r=-0{,}01$, IC incluye 0. \emph{Fuente: elaboración propia.}}
\end{column}
\end{columns}
\end{frame}

\begin{frame}{Resultado 4: el desbalance no gobierna la estabilidad}
\begin{columns}[T]
\begin{column}{0.5\textwidth}
\begin{itemize}\small
  \item Esperábamos degradación al agravarse el desbalance. \textbf{No la hay.}
  \item Rango estrecho: de $0{,}696$ (1:1) a $0{,}765$ (1:50), amplitud de \textbf{siete centésimas}, menor que las once que separan a las arquitecturas.
  \item \textbf{Kruskal-Wallis $p = 0{,}18$}: no se rechaza la igualdad entre los cinco escenarios. Tamaño de efecto $\eta^2 = 0{,}05$, pequeño frente al $0{,}13$ de la arquitectura.
  \item El valor \textbf{más bajo} es 1:1, el más equilibrado: la variación no tiene dirección.
  \item El \textbf{balanceo} tampoco pesa: $\eta^2 \le 0{,}01$ en las tres dimensiones.
\end{itemize}
\vspace{0.4em}
\begin{block}{}
\footnotesize Implicación práctica: el balanceo puede elegirse por \textbf{rendimiento}, sin temer que degrade la interpretabilidad.
\end{block}
\end{column}
\begin{column}{0.5\textwidth}\centering
\figcard{chapter_4/images_ch4/estabilidad_escenario.png}
\figcap{Estabilidad media por escenario de desbalance sobre Elliptic, promediada sobre las 3 semillas de modelo. \emph{Fuente: elaboración propia.}}
\end{column}
\end{columns}
\end{frame}

\begin{frame}{Rendimiento y colapso validación $\rightarrow$ test}
\begin{columns}[c]
\begin{column}{0.60\textwidth}
\centering
\figcard[\linewidth]{colapso_val_test.png}
\figcap{Rendimiento en validación y test sobre los mismos modelos: PR-AUC y precisión@50 colapsan, mientras que el ROC-AUC cae mucho menos y aparenta un modelo aceptable. \emph{Fuente: elaboración propia.}}
\end{column}
\begin{column}{0.40\textwidth}
\small
\begin{itemize}
\item El modelo \textbf{aprende en validación} y \textbf{colapsa en test}: es el \emph{shift} temporal del dataset (los patrones de lavado cambian con el tiempo). Propiedad del dato, no del método.
\item El \textbf{ROC-AUC engaña} bajo desbalance ($0{,}88$): tomamos \textbf{PR-AUC} y \textbf{precisión@k} como primarias.
\item La estabilidad se estudia sobre \textbf{verdaderos positivos de validación} (encuadre declarado).
\end{itemize}
\end{column}
\end{columns}
\end{frame}

\begin{frame}{Rigor métrico: por qué PR-AUC y no ROC-AUC}
\begin{center}
\small
\tabcap{Métricas de ordenamiento en validación y test (mismos modelos).}
\renewcommand{\arraystretch}{1.3}
\begin{tabular}{lcc}
\toprule
\textbf{Métrica de ordenamiento} & \textbf{Validación} & \textbf{Test} \\
\midrule
ROC-AUC & 0,884 & 0,653 \\
\textbf{\color{amber}PR-AUC} & \textbf{0,367} & \textbf{0,017} \\
Precisión en los primeros 50 & 0,657 & 0,020 \\
Precisión en los primeros 100 & 0,640 & 0,022 \\
\bottomrule
\end{tabular}
\end{center}
\vspace{0.5em}
{\footnotesize El ROC-AUC de 0,884 sugeriría un clasificador casi excelente. El PR-AUC de 0,367 sobre los mismos modelos revela la dificultad real: bajo desbalance extremo, el eje de falsos positivos queda dominado por la clase mayoritaria y permanece bajo aunque el modelo no distinga la clase rara.}
\end{frame}

\begin{frame}{Las tres hipótesis, y qué pasó con cada una}
\begin{center}
\small
\tabcap{Hipótesis de trabajo y su veredicto empírico.}
\renewcommand{\arraystretch}{1.5}
\begin{tabular}{p{0.42\textwidth}p{0.15\textwidth}p{0.32\textwidth}}
\toprule
\textbf{Lo que predijimos} & \textbf{Veredicto} & \textbf{La evidencia} \\
\midrule
\textbf{H1.} La estabilidad se degradaría con el desbalance. & \textcolor{coral}{\textbf{Refutada}} & Variación sin dirección: $p = 0{,}18$, mínimo en 1:1 (el más equilibrado). \\
\textbf{H2.} TAGCN sería la más estable (alcance multi-hop). & \textcolor{coral}{\textbf{Refutada}} & TAGCN en el grupo bajo en las 3 semillas. \\
\textbf{H3.} Un explicador más estable señalaría mejor el patrón. & \textcolor{coral}{\textbf{Refutada}} & Puente nulo: $r=-0{,}01$, IC incluye el cero. \\
\bottomrule
\end{tabular}
\end{center}
\vspace{0.6em}
{\footnotesize Las tres se cayeron y las tres se reportan. Haber escrito las predicciones antes de ver los datos es lo que convierte su refutación en un resultado.}
\end{frame}

% ##################################################################### 5
\section{Conclusiones}

\begin{frame}{Conclusiones del proyecto}
\begin{enumerate}\small
  \item \textbf{\color{amber}Tres dimensiones independientes} (estabilidad, plausibilidad, fidelidad): no existe una única ``buena explicación'' y el mejor explicador depende del objetivo de la auditoría.
  \item \textbf{\color{amber}La estabilidad por arquitectura son dos grupos} (alto con GAT y GCN, bajo con GraphSAGE y TAGCN), y la partición \textbf{concuerda} entre datos reales y sintéticos.
  \item \textbf{\color{amber}El explicador es la palanca dominante} ($\eta^2$ de $0{,}36$ a $0{,}64$). El balanceo es despreciable ($\eta^2 \le 0{,}01$) y el escenario de desbalance, pequeño y no significativo ($\eta^2 \le 0{,}05$).
  \item \textbf{\color{amber}El más plausible no es el más fiel} (PGExplainer vs GNNExplainer): elegir según se quiera \emph{reconocer el patrón} o \emph{auditar el modelo}.
  \item \textbf{\color{amber}Aporte metodológico}: dos artefactos de medición corregidos + distinguir reproducibilidad de \emph{pesos} vs de \emph{conclusiones}.
  \item \textbf{\color{amber}Las tres hipótesis se refutaron} y se reportan tal cual, sin ajustar el análisis para salvarlas.
\end{enumerate}
\end{frame}

\begin{frame}{Los cuatro objetivos, respondidos}
\begin{center}
\footnotesize
\tabcap{Respuesta a los objetivos de investigación.}
\renewcommand{\arraystretch}{1.5}
\begin{tabular}{p{0.28\textwidth}p{0.62\textwidth}}
\toprule
\textbf{Objetivo} & \textbf{Respuesta} \\
\midrule
\textbf{O1.} Desbalance vs robustez explicativa & No es el factor dominante. Variación sin dirección entre escenarios ($p = 0{,}18$) y balanceo despreciable. La palanca real es el \textbf{explicador}. \\
\textbf{O2.} Resiliencia de las 4 arquitecturas & No hay cuatro posiciones, hay \textbf{dos grupos}. Se replica en ambos regímenes de densidad. \\
\textbf{O3.} Estrategias de balanceo & Despreciable en las tres dimensiones ($\eta^2 \le 0{,}01$). Elegible por rendimiento. \\
\textbf{O4.} Matriz de recomendación & No existe tríada óptima única. Recomendación \textbf{condicional al propósito}. \\
\bottomrule
\end{tabular}
\end{center}
\end{frame}

\begin{frame}{Matriz de recomendación}
\begin{center}
\tabcap{Matriz de recomendación por objetivo de auditoría.}
\renewcommand{\arraystretch}{1.4}
\begin{tabular}{ll}
\toprule
\textbf{Objetivo operativo} & \textbf{Configuración} \\
\midrule
Auditabilidad / estabilidad & \textcolor{teal}{\textbf{grupo alto: GAT o GCN}} \\
Recuperar el patrón (plausibilidad) & \textcolor{seafoam}{\textbf{PGExplainer}} \\
Fidelidad al modelo & \textcolor{amber}{\textbf{GNNExplainer}} \\
Estabilidad interna del método & \textcolor{coral}{\textbf{GNNShap}} \\
\bottomrule
\end{tabular}
\end{center}
\vspace{0.6em}
{\footnotesize\color{muted} La recomendación de arquitectura es \textbf{de grupo}, no de una arquitectura concreta: dentro del grupo alto la evidencia no distingue, así que puede atender a coste o disponibilidad.}
\end{frame}

\begin{frame}{Contribuciones y limitaciones}
\begin{columns}[T]
\begin{column}{0.5\textwidth}
\begin{block}{Contribuciones}
\begin{itemize}\small
  \item Tres dimensiones independientes.
  \item Dos artefactos + dos \emph{bugs} de PGExplainer.
  \item Generador sintético con \emph{ground-truth} por arista.
  \item Reproducibilidad de \textbf{pesos} vs de \textbf{conclusiones}.
  \item Matriz de recomendación por objetivo.
\end{itemize}
\end{block}
\end{column}
\begin{column}{0.5\textwidth}
\begin{alertblock}{Limitaciones (declaradas)}
\begin{itemize}\small
  \item La inferencia fuerte proviene del eje sintético.
  \item El clasificador colapsa en test (\emph{shift} temporal).
  \item Un solo \emph{dataset} real, anonimizado.
  \item Con 3 semillas se sostiene el \textbf{grupo}, no el orden fino dentro.
\end{itemize}
\end{alertblock}
\end{column}
\end{columns}
\end{frame}

\begin{frame}[standout]
Gracias \textperiodcentered\ Preguntas
\vspace{0.5em}

{\small\color{greytx} Estabilidad $\neq$ plausibilidad $\neq$ fidelidad: el mejor explicador depende del objetivo.\\
Trabajo futuro: Elliptic2, AMLSim, GNNs temporales, GraphSMOTE.}
\vspace{0.5em}

{\footnotesize\color{muted} Alejandro Gómez Huertas \textperiodcentered\ Juan Diego Garzón Ovalle}
\end{frame}

% ##################################################################### RESPALDO
\appendix

\begin{frame}[allowframebreaks]{Referencias seleccionadas}
\scriptsize
\begin{itemize}
  \item Weber M y colaboradores (2019). \emph{Anti-Money Laundering in Bitcoin: Experimenting with Graph Convolutional Networks for Financial Forensics.} arXiv:1908.02591.
  \item Kipf T y Welling M (2017). \emph{Semi-Supervised Classification with Graph Convolutional Networks.} ICLR.
  \item Hamilton W, Ying R y Leskovec J (2017). \emph{Inductive Representation Learning on Large Graphs.} NeurIPS.
  \item Veličković P y colaboradores (2018). \emph{Graph Attention Networks.} ICLR.
  \item Du J y colaboradores (2017). \emph{Topology Adaptive Graph Convolutional Networks.} arXiv:1710.10370.
  \item Ying R y colaboradores (2019). \emph{GNNExplainer: Generating Explanations for Graph Neural Networks.} NeurIPS.
  \item Luo D y colaboradores (2020). \emph{Parameterized Explainer for Graph Neural Network.} NeurIPS.
  \item Akkas S y Azad A (2024). \emph{GNNShap: Scalable and Accurate GNN Explanation using Shapley Values.} WWW.
  \item Lundberg S y Lee S (2017). \emph{A Unified Approach to Interpreting Model Predictions.} NeurIPS.
  \item Kosan M y colaboradores (2023). \emph{GNNX-BENCH: Unravelling the Utility of Perturbation-based GNN Explainers through In-depth Benchmarking.} arXiv:2310.01794.
  \item Agarwal C, Zitnik M y Lakkaraju H (2022). \emph{Probing GNN Explainers: A Rigorous Theoretical and Empirical Analysis of GNN Explanation Methods.} AISTATS.
  \item Lawal O, Okolie A y Obunadike C (2025). \emph{An Explainable Graph Neural Network Framework for Anti-Money Laundering in Cryptocurrency Transactions Using the Elliptic Dataset.} IJNSA 17(6). DOI: 10.5121/ijnsa.2025.17602.
  \item He X y colaboradores (2026). \emph{An Explainable Graph Neural Network Framework for Illicit Financial Transaction Detection.} Applied Intelligence 56. DOI: 10.1007/s10489-026-07138-9.
  \item Zhao T, Zhang X y Wang S (2021). \emph{GraphSMOTE: Imbalanced Node Classification on Graphs with Graph Neural Networks.} WSDM.
  \item Lin T y colaboradores (2017). \emph{Focal Loss for Dense Object Detection.} ICCV.
\end{itemize}
\end{frame}

\begin{frame}{Respaldo \textperiodcentered\ disociación: valores exactos}
\begin{center}
\small
\tabcap{Plausibilidad y Fidelity+ por explicador (eje sintético, g0, 3 semillas). Cada explicador se compara contra la línea base de azar de su propia métrica.}
\renewcommand{\arraystretch}{1.3}
\begin{tabular}{lccc}
\toprule
\textbf{Explicador} & \textbf{Plaus. aristas} & \textbf{Plaus. features} & \textbf{Fidelity+} \\
\midrule
GNNExplainer & 0,50 & \textbf{0,45} & \textbf{0,56} \\
PGExplainer & \textbf{0,80} & n/d & 0,11 \\
GNNShap & n/d & 0,15 & 0,46 \\
\midrule
\emph{Azar (línea base)} & \emph{0,40} & \emph{0,08} & no aplica \\
\bottomrule
\end{tabular}
\end{center}
\vspace{0.4em}
{\footnotesize Las celdas n/d corresponden a métricas que el explicador no produce por diseño: PGExplainer no ordena \emph{features} y GNNShap no genera máscara de aristas. Contra su propia línea base, PGExplainer duplica el azar en aristas y GNNExplainer lo sextuplica en \emph{features}. PGExplainer supera a GNNExplainer en plausibilidad de aristas: Wilcoxon pareado $p\approx 2{,}6\times10^{-35}$ ($n=225$ parejas). \textbf{GNNShap} es el más estable internamente entre ejecuciones ($0{,}98$), aunque no lidere plausibilidad ni fidelidad.}
\end{frame}

\begin{frame}{Respaldo \textperiodcentered\ estabilidad por semilla de modelo}
\begin{center}
\small
\tabcap{Estabilidad por semilla de modelo (Elliptic, GNNExplainer).}
\renewcommand{\arraystretch}{1.25}
\begin{tabular}{lccccc}
\toprule
\textbf{Arquitectura} & \textbf{s42} & \textbf{s43} & \textbf{s44} & \textbf{Media} & \textbf{IC95\,\%} \\
\midrule
GAT       & 0,766 & 0,789 & 0,790 & \textbf{0,782} & 0,758 a 0,805 \\
GCN       & 0,783 & 0,733 & 0,757 & \textbf{0,758} & 0,724 a 0,787 \\
GraphSAGE & 0,756 & 0,713 & 0,738 & 0,735 & 0,710 a 0,756 \\
TAGCN     & 0,586 & 0,693 & 0,736 & 0,672 & 0,632 a 0,717 \\
\bottomrule
\end{tabular}
\end{center}
\vspace{0.5em}
{\footnotesize \textbf{GAT y GCN se permutan} entre semillas: por eso no nombramos un ganador único. \textbf{TAGCN} tiene la mayor dispersión (dt 0,077, el triple): su lugar en el grupo bajo es firme, su valor central admite margen.}
\end{frame}

\begin{frame}{Respaldo \textperiodcentered\ detalle estadístico de la partición}
\begin{center}
\footnotesize
\tabcap{Wilcoxon pareado por pares de arquitecturas (Elliptic y sintético).}
\renewcommand{\arraystretch}{1.2}
\begin{tabular}{lcc}
\toprule
\textbf{Wilcoxon pareado} & \textbf{Elliptic} & \textbf{Sintético} \\
\midrule
GAT vs GCN \emph{(grupo alto)} & 0,165 & 0,0013 \\
GraphSAGE vs TAGCN \emph{(grupo bajo)} & 0,096 & 0,519 \\
\midrule
GAT vs GraphSAGE & 0,0033 & $<10^{-4}$ \\
GAT vs TAGCN & $<10^{-4}$ & $<10^{-4}$ \\
GCN vs GraphSAGE & 0,047 & $<10^{-4}$ \\
GCN vs TAGCN & 0,0067 & $<10^{-4}$ \\
\bottomrule
\end{tabular}
\end{center}
\vspace{0.5em}
{\footnotesize La única diferencia intragrupo significativa es GCN sobre GAT en el sintético, por un margen de \textbf{0,006} sin relevancia práctica, detectable solo por el mayor $n$ de ese eje.}
\end{frame}

\begin{frame}{Respaldo \textperiodcentered\ ¿por qué GNNExplainer para el ranking?}
\begin{center}
\small
\tabcap{Estabilidad (Spearman) por explicador y arquitectura en Elliptic.}
\renewcommand{\arraystretch}{1.25}
\begin{tabular}{lcccc}
\toprule
\textbf{Explicador (Elliptic)} & \textbf{GAT} & \textbf{GCN} & \textbf{GraphSAGE} & \textbf{TAGCN} \\
\midrule
\textbf{\color{teal}GNNExplainer} & 0,78 & 0,76 & 0,74 & 0,67 \\
PGExplainer & 0,00 & 0,00 & 0,00 & 0,00 \\
GNNShap & 0,97$^{*}$ & 0,97 & 0,95 & 0,97$^{*}$ \\
\bottomrule
\end{tabular}
\end{center}
\vspace{0.5em}
{\footnotesize \textbf{\color{teal}GNNExplainer} es el único que discrimina entre arquitecturas: PGExplainer degenera (Spearman nulo) y GNNShap se satura cerca de 0,96, indistinguible. Además es el explicador común a los dos ejes, lo que hace comparable la partición. GNNExplainer, a 3 semillas en las cuatro arquitecturas. GNNShap, a 3 semillas en GCN y GraphSAGE. Las marcadas con $^{*}$, GAT y TAGCN, solo en la semilla 42.}
\end{frame}

\begin{frame}{Respaldo \textperiodcentered\ curvas ROC y precisión-recall}
\begin{center}
\figcard[0.96\linewidth]{curvas_pr_roc.png}
\end{center}
\figcap{Curvas ROC (izquierda) y de precisión y exhaustividad (derecha) sobre validación y test, mismos modelos. El ROC cae de 0,88 a 0,65, mientras que la precisión y exhaustividad se desploma casi a la tasa base de transacciones ilícitas: por eso la métrica primaria es PR-AUC y no ROC-AUC. \emph{Fuente: elaboración propia.}}
\end{frame}

\end{document}
